\begin{textsample}{POS Dim 1 – 2015 – Score 22.00 – 2015/t001608.txt}  \label{ex:f1_pos_040}
Well, you know, sometimes \textbf{the} most important things come in \textbf{the} smallest packages.
I am going to try to convince you, in \textbf{the} 15 minutes I have, that microbes have a lot to say about questions such as,"Are we alone?"and they can tell us more about not only life in our solar system but also maybe beyond, and this is why I am tracking them down in \textbf{the} most impossible places on Earth, in extreme environments where conditions are really pushing them to \textbf{the} brink of survival.
Actually, sometimes me too, when I ' m trying to follow them too close.
But here ' s \textbf{the} thing : We are \textbf{the} only advanced civilization in \textbf{the} solar system, but that doesn ' t mean that there is no microbial life nearby.
In fact, \textbf{the} \textbf{planets} and moons you see here could host life — all of them — and we know that, and it ' s a strong possibility.
And if we were going to find life on those moons and \textbf{planets}, then we would answer questions such as, are we alone in \textbf{the} solar system?
Where are we coming from?
Do we have family in \textbf{the} neighborhood?
Is there life beyond our solar system?
And we can ask all those questions because there has been a revolution in our understanding of what a habitable \textbf{planet} is, and today, a habitable \textbf{planet} is a \textbf{planet} that has a zone where water can stay stable, but to me this is a \textbf{horizontal} definition of habitability, because it involves a distance to a star, but there is another \textbf{dimension} to habitability, and this is a vertical \textbf{dimension}.
Think of it as conditions in \textbf{the} subsurface of a \textbf{planet} where you are very far away from a sun, but you still have water, energy, nutrients, which for some of them means food, and a protection.
And when you look at \textbf{the} Earth, very far away from any sunlight, deep in \textbf{the} ocean, you have life thriving and it uses only \textbf{chemistry} for life processes.
So when you think of it at that point, all walls collapse.
You have no limitations, basically.
And if you have been looking at \textbf{the} headlines lately, then you will see that we have discovered a subsurface ocean on Europa, on Ganymede, on Enceladus, on Titan, and now we are finding a geyser and hot \textbf{springs} on Enceladus, Our solar system is turning into a giant spa.
For anybody who has gone to a spa knows how much microbes like that, right?
So at that point, think also about Mars.
There is no life possible at \textbf{the} \textbf{surface} of Mars today, but it might still be hiding \textbf{underground}.
So, we have been making progress in our understanding of habitability, but we also have been making progress in our understanding of what \textbf{the} signatures of life are on Earth.
And you can have what we call \textbf{organic} molecules, and these are \textbf{the} bricks of life, and you can have \textbf{fossils}, and you can minerals, biominerals, which is due to \textbf{the} reaction between bacteria and \textbf{rocks}, and of course you can have gases in \textbf{the} atmosphere.
And when you look at those tiny green algae on \textbf{the} right of \textbf{the} \textbf{slide} here, they are \textbf{the} direct descendants of those who have been pumping oxygen a billion years ago in \textbf{the} atmosphere of \textbf{the} Earth.
When they did that, they poisoned 90 percent of \textbf{the} life at \textbf{the} \textbf{surface} of \textbf{the} Earth, but they are \textbf{the} reason why you are breathing this air today.
But as much as our understanding grows of all of these things, there is one question we still cannot answer, and this is, where are we coming from?
And you know, it ' s getting worse, because we won ' t be able to find \textbf{the} physical evidence of where we are coming from on this \textbf{planet}, and \textbf{the} reason being is that anything that is older than four billion years is gone.
All record is gone, erased by \textbf{plate} tectonics and erosion.
This is what I call \textbf{the} Earth ' s biological horizon.
Beyond this horizon we don ' t know where we are coming from.
So is everything lost?
Well, maybe not.
And we might be able to find evidence of our own origin in \textbf{the} most unlikely place, and this place in Mars.
How is this possible?
Well clearly at \textbf{the} beginning of \textbf{the} solar system, Mars and \textbf{the} Earth were bombarded by giant \textbf{asteroids} and comets, and there were ejecta from these impacts all over \textbf{the} place.
Earth and Mars kept throwing \textbf{rocks} at each other for a very long time.
Pieces of \textbf{rocks} landed on \textbf{the} Earth.
Pieces of \textbf{the} Earth landed on Mars.
So clearly, those two \textbf{planets} may have been seeded by \textbf{the} same material.
So yeah, maybe Granddady is sitting there on \textbf{the} \textbf{surface} and waiting for us.
But that also means that we can go to Mars and try to find traces of our own origin.
Mars may hold that secret for us.
This is why Mars is so special to us.
But for that to happen, Mars needed to be habitable at \textbf{the} time when conditions were right.
So was Mars habitable?
We have a number of missions telling us exactly \textbf{the} same thing today.
At \textbf{the} time when life appeared on \textbf{the} Earth, Mars did have an ocean, it had volcanoes, it had lakes, and it had deltas like \textbf{the} beautiful picture you see here.
This picture was sent by \textbf{the} Curiosity rover only a few weeks ago.
It shows \textbf{the} remnants of a delta, and this picture tells us something : water was abundant and stayed founting at \textbf{the} \textbf{surface} for a very long time.
This is good news for life.
Life \textbf{chemistry} takes a long time to actually happen.
So this is extremely good news, but does that mean that if we go there, life will be easy to find on Mars?
Not necessarily.
Here ' s what happened : At \textbf{the} time when life exploded at \textbf{the} \textbf{surface} of \textbf{the} Earth, then everything went south for Mars, literally. \textbf{The} atmosphere was stripped away by solar winds, Mars lost its magnetosphere, and then cosmic rays and U.
V. bombarded \textbf{the} \textbf{surface} and water escaped to space and went \textbf{underground}.
So if we want to be able to understand, if we want to be able to find those traces of \textbf{the} signatures of life at \textbf{the} \textbf{surface} of Mars, if they are there, we need to understand what was \textbf{the} impact of each of these events on \textbf{the} preservation of its record.
Only then will we be able to know where those signatures are hiding, and only then will we be able to send our rover to \textbf{the} right places where we can sample those \textbf{rocks} that may be telling us something really important about who we are, or, if not, maybe telling us that somewhere, independently, life has appeared on another \textbf{planet}.
So to do that, it ' s easy.
You only need to go back 3. 5 billion years ago in \textbf{the} past of a \textbf{planet}.
We just need a time machine.
Easy, right?
Well, actually, it is.
Look around you — that ' s \textbf{planet} Earth.
This is our time machine.
Geologists are using it to go back in \textbf{the} past of our own \textbf{planet}.
I am using it a little bit differently.
I use \textbf{planet} Earth to go in very extreme environments where conditions were similar to those of Mars at \textbf{the} time when \textbf{the} climate changed, and there I ' m trying to understand what happened.
What are \textbf{the} signatures of life?
What is left?
How are we going to find it?
So for one moment now I ' m going to take you with me on a trip into that time machine.
And now, what you see here, we are at 4, 500 meters in \textbf{the} Andes, but in fact we are less than a billion years after \textbf{the} Earth and Mars formed. \textbf{The} Earth and Mars will have looked pretty much exactly like that — volcanoes everywhere, evaporating lakes everywhere, minerals, hot \textbf{springs}, and then you see those mounds on \textbf{the} shore of those lakes?
Those are built by \textbf{the} descendants of \textbf{the} first organisms that gave us \textbf{the} first \textbf{fossil} on Earth.
But if we want to understand what ' s going on, we need to go a little further.
And \textbf{the} other thing about those sites is that exactly like on Mars three and a half billion years ago, \textbf{the} climate is changing very fast, and water and ice are \textbf{disappearing}.
But we need to go back to that time when everything changed on Mars, and to do that, we need to go higher.
Why is that?
Because when you go higher, \textbf{the} atmosphere is getting thinner, it ' s getting more unstable, \textbf{the} temperature is getting cooler, and you have a lot more U.
V. radiation.
Basically, you are getting to those conditions on Mars when everything changed.
So I was not promising anything about a leisurely trip on \textbf{the} time machine.
You are not going to be sitting in that time machine.
You have to haul 1, 000 \textbf{pounds} of equipment to \textbf{the} summit of this 20, 000-foot volcano in \textbf{the} Andes here.
That ' s about 6, 000 meters.
And you also have to sleep on 42-degree \textbf{slopes} and really hope that there won ' t be any \textbf{earthquake} that night.
But when we get to \textbf{the} summit, we actually find \textbf{the} lake we came for.
At this \textbf{altitude}, this lake is experiencing exactly \textbf{the} same conditions as those on Mars three and a half billion years ago.
And now we have to change our voyage into an inner voyage inside that lake, and to do that, we have to remove our mountain gear and actually don suits and go for it.
But at \textbf{the} time we enter that lake, at \textbf{the} very moment we enter that lake, we are stepping back three and a half billion years in \textbf{the} past of another \textbf{planet}, and then we are going to get \textbf{the} answer came for.
Life is everywhere, absolutely everywhere.
Everything you see in this picture is a living organism.
Maybe not so \textbf{the} diver, but everything else.
But this picture is very deceiving.
Life is abundant in those lakes, but like in many places on Earth right now and due to climate change, there is a huge loss in biodiversity.
In \textbf{the} samples that we took back home, 36 percent of \textbf{the} bacteria in those lakes were composed of three species, and those three species are \textbf{the} ones that have survived so far.
Here ' s another lake, right next to \textbf{the} first one. \textbf{The} red color you see here is not due to minerals.
It ' s actually due to \textbf{the} presence of a tiny algae.
In this region, \textbf{the} U.
V. radiation is really nasty.
Anywhere on Earth, 11 is considered to be extreme.
During U.
V. storms there, \textbf{the} U.
V.
Index reaches 43.
SPF 30 is not going to do anything to you over there, and \textbf{the} water is so transparent in those lakes that \textbf{the} algae has nowhere to hide, really, and so they are developing their own sunscreen, and this is \textbf{the} red color you see.
But they can adapt only so far, and then when all \textbf{the} water is gone from \textbf{the} \textbf{surface}, microbes have only one solution \textbf{left} : They go \textbf{underground}.
And those microbes, \textbf{the} \textbf{rocks} you see in that \textbf{slide} here, well, they are actually living inside \textbf{rocks} and they are using \textbf{the} protection of \textbf{the} translucence of \textbf{the} \textbf{rocks} to get \textbf{the} good part of \textbf{the} U.
V. and \textbf{discard} \textbf{the} part that could actually damage their DNA.
And this is why we are taking our rover to train them to \textbf{search} for life on Mars in these areas, because if there was life on Mars three and a half billion years ago, it had to use \textbf{the} same strategy to actually protect itself.
Now, it is pretty obvious that going to extreme environments is helping us very much for \textbf{the} exploration of Mars and to prepare missions.
So far, it has helped us to understand \textbf{the} geology of Mars.
It has helped to understand \textbf{the} past climate of Mars and its \textbf{evolution}, but also its habitability potential.
Our most recent rover on Mars has discovered traces of \textbf{organics}.
Yeah, there are \textbf{organics} at \textbf{the} \textbf{surface} of Mars.
And it also discovered traces of methane.
And we don ' t know yet if \textbf{the} methane in question is really from geology or \textbf{biology}.
Regardless, what we know is that because of \textbf{the} discovery, \textbf{the} hypothesis that there is still life present on Mars today remains a viable one.
So by now, I think I have convinced you that Mars is very special to us, but it would be a mistake to think that Mars is \textbf{the} only place in \textbf{the} solar system that is interesting to find potential microbial life.
And \textbf{the} reason is because Mars and \textbf{the} Earth could have a common root to their tree of life, but when you go beyond Mars, it ' s not that easy.
Celestial mechanics is not making it so easy for an exchange of material between \textbf{planets}, and so if we were to discover life on those \textbf{planets}, it would be different from us.
It would be a different type of life.
But in \textbf{the} end, it might be just us, it might be us and Mars, or it can be many trees of life in \textbf{the} solar system.
I don ' t know \textbf{the} answer yet, but I can tell you something : No matter what \textbf{the} result is, no matter what that magic number is, it is going to give us a standard by which we are going to be able to measure \textbf{the} life potential, abundance and diversity beyond our own solar system.
And this can be achieved by our generation.
This can be our legacy, but only if we dare to explore.
Now, finally, if somebody tells you that looking for alien microbes is not cool because you cannot have a philosophical conversation with them, let me show you why and how you can tell them they ' re wrong.
Well, \textbf{organic} material is going to tell you about environment, about complexity and about diversity.
DNA, or any information carrier, is going to tell you about adaptation, about \textbf{evolution}, about survival, about planetary changes and about \textbf{the} transfer of information.
All together, they are telling us what started as a microbial pathway, and why what started as a microbial pathway sometimes ends up as a civilization or sometimes ends up as a dead end.
Look at \textbf{the} solar system, and look at \textbf{the} Earth.
On Earth, there are many intelligent species, but only one has achieved technology.
Right here in \textbf{the} journey of our own solar system, there is a very, very powerful message that says here ' s how we should look for alien life, small and big.
So yeah, microbes are talking and we are listening, and they are taking us, one \textbf{planet} at a time and one moon at a time, towards their big brothers out there.
And they are telling us about diversity, they are telling us about abundance of life, and they are telling us how this life has survived thus far to reach civilization, intelligence, technology and, indeed, philosophy.
Thank you.

% matched lemmas: altitude, asteroid, biology, chemistry, dimension, disappear, discard, earthquake, evolution, fossil, horizontal, left, organic, planet, plate, pound, rock, search, slide, slope, spring, surface, the, underground
\end{textsample}
